\documentclass[11pt]{article}
\usepackage[margin=1in]{geometry}
\usepackage{amsmath,amssymb}
\usepackage{booktabs}
\usepackage{microtype}
\usepackage[hidelinks]{hyperref}

\title{\bf When Is a Market Estimate Trustworthy?\\
\large Error Bounds for Margin-Constrained Market Disaggregation}
\author{}
\date{}

\begin{document}
\maketitle

\begin{abstract}
Most market estimates are built by taking a published total and multiplying it
down through geography, industry, and segment shares. The method assumes spend is
a measure over a partitionable universe, so that any market is a subset sum. That
assumption is sound. The problem is that it says nothing about whether a
particular estimate is any good, and practitioners need exactly that.

We show the question has a clean answer. Iterative proportional fitting (IPF)
reproduces its input margins exactly, so its residual is orthogonal to the span
of the observed constraints. Only the part of a query pointing outside that span
can be wrong. This gives a reliability statistic computable before you answer,
with no access to ground truth. Tested against real published joint
distributions from the U.S. County Business Patterns (40,560 cells, 6.1M
establishments), it is calibrated to within 1.15$\times$ across six seeds and cuts
the dispersion of prediction error by about a third against a query-size
baseline. It does \emph{not} rank queries better than query size alone, and we
report that head-to-head in full because the temptation to claim otherwise is
strong and the fair comparison is easy to get wrong.

We then ask how much better modelling could do, and the answer is: not much.
Three independent approaches, a parametric interaction term, low-rank
completion, and covariate regression, each recover somewhere between 2\% and
11\% of the residual. Roughly nine tenths of it is cell-specific and out of reach
of anything transferable. Retaining past estimates as constraints helps by 87\% on queries near what was
already observed and by an amount not distinguishable from zero on queries that
are not, and the benefit saturates after a \emph{single} nearby observation. The
value of such a store therefore scales with the number of distinct neighbourhoods
it covers, not with how many records it holds. Moving from
establishment counts to dollars costs a steady factor of about 1.5$\times$, which
is a tax rather than a wall. We also document two failures that leave no trace in
the output: IPF converges silently and order-dependently when its margins
disagree, and disclosure control removes cells in a pattern built specifically to
defeat inference from margins.
\end{abstract}

\section{Introduction}

A market estimate usually starts with a number someone else published and works
downward. Global category total, times a country share, times a segment share,
times a channel share. The logic holds up: expenditure is an accounting
identity, so there is a finite pool, and any market is a piece of it. Treat spend
as a measure over disjoint cells and market size becomes a subset sum, with
disaggregation reduced to interpolation under a known constraint.

Two things get in the way. Markets are not cells of the census lattice, because a
market is really an equivalence class under substitutability, and that cuts
diagonally across geography and industry codes. More seriously, the joint
distribution over cells is not published. What gets published is a handful of
low-dimensional projections of it. Recovering a joint from its margins is badly
underdetermined, and the usual remedy, maximum-entropy fill-in by IPF, picks one
point out of a large feasible set rather than finding the right one.

So a disaggregated estimate can be excellent or worthless, and nothing about the
output tells you which. That is the gap this paper closes.

\paragraph{What is here.} We restate markets as linear functionals over spend
atoms, which lets overlapping markets coexist without breaking additivity. We
derive a per-query error identity and turn it into a calibrated statistic. We
measure realised error against published joints rather than simulated ones. We
put a ceiling on what interaction modelling can add. And we document two failure
modes that current practice does not surface.

The mathematics is not new. The orthogonality in Section~\ref{sec:identity} is
elementary linear algebra. IPF goes back to Deming and Stephan~\cite{deming1940},
its convergence theory to Csisz\'ar~\cite{csiszar1975}, and estimation for small
domains has a mature literature of its own~\cite{rao2015}. What we offer is the
packaging of an elementary fact into a decision rule, plus the calibration that
makes it usable.

\section{Setup}

Let $\mathcal{A}$ be a finite set of \emph{atoms}: the finest cells at which
spend is defined, indexed by geography, industry, size band, and period. Atoms
are disjoint and exhaustive by construction. Let
$T \in \mathbb{R}_{\ge 0}^{\mathcal{A}}$ hold the true spend.

A \emph{market} is not a cell. It is a linear functional. A query
$q \in [0,1]^{\mathcal{A}}$ gives each atom a membership weight, and the market's
size is $\langle q, T\rangle$. Fractional weights carry real content: whether
convenience-store coffee belongs to the coffee-shop market is a judgement about
substitution, and this is where it lives. The representation also lets two
markets overlap freely, since overlap just means shared atoms rather than
colliding sets. That resolves an awkward fact about published figures, which is
that category TAMs across an economy sum to far more than its GDP.

An \emph{observation} is a pair $(a,y)$ with $\langle a, T\rangle = y$. Published
marginal tables are observations whose incidence vectors are indicators of
fibres. Write $V = \mathrm{span}\{a_1,\dots,a_m\}$ for the constraint span.

Given observations, IPF returns the $H$ minimising
$D_{\mathrm{KL}}(H \Vert \mathbf{1})$ subject to
$\langle a_i, H\rangle = y_i$. It converges geometrically, being a contraction in
the Hilbert projective metric~\cite{sinkhornrate}, and is the same algorithm as
Sinkhorn--Knopp. When the constraint set is \emph{decomposable}, the
maximum-likelihood solution has a closed form and no iteration is
needed~\cite{decomposable}, which is worth checking before running a loop.

\section{An exact error identity}
\label{sec:identity}

IPF reproduces its inputs exactly: $\langle a_i, H\rangle = \langle a_i, T\rangle$
for every observed $a_i$. So the residual $R = H - T$ satisfies
$\langle a_i, R\rangle = 0$ for all $i$, which is to say $R \perp V$ under the
unweighted inner product. For any query $q$, writing $P_V$ for orthogonal
projection onto $V$ and $q_\perp = q - P_V(q)$,
\begin{equation}
  \mathrm{err}(q) \;=\; \langle q, R\rangle \;=\; \langle q_\perp, R\rangle .
  \label{eq:identity}
\end{equation}

Only the part of a query pointing outside what you measured can be wrong. Two
things follow.

\paragraph{Some queries are exact.} If $q \in V$ then $\mathrm{err}(q) = 0$
identically. Concretely, a query that pins some axes and integrates another away
is answered exactly whenever the matching margin was observed. Asking for every
construction firm in a set of counties, at any size, is free.

\paragraph{A statistic you can compute first.} Cauchy--Schwarz gives
$|\mathrm{err}(q)| \le \lVert q_\perp \rVert \, \lVert R \rVert$. The quantity
that matters is the \emph{unnormalised} norm $\lVert q_\perp \rVert$, not the
alignment ratio $\lVert P_V q\rVert^2/\lVert q\rVert^2$, which divides out
precisely the scale carrying the signal. We learned this the hard way: the ratio
turned out to have essentially no predictive power. Treating $R$ as isotropic to
first order, the predicted relative error is
\begin{equation}
  \widehat{\varepsilon}(q) \;=\; \kappa \cdot
  \frac{\lVert q - P_V(q) \rVert_2}{\langle q, H\rangle},
  \label{eq:predictor}
\end{equation}
with $\kappa$ a single scalar calibrated once per lattice. Everything on the
right is available without ground truth. The query is given, $\langle q,H\rangle$
is the estimate itself, and $P_V(q)$ comes from the constraint structure by
backfitting. The statistic returns exactly zero on queries in $V$ and shrinks as
queries grow, through the denominator, so it folds two separately visible effects
into one number.

\section{Data and protocol}

Synthetic data would let us choose the interaction strength that decides the
answer, which defeats the purpose. So we used a published joint instead. The U.S.
Census County Business Patterns (CBP) county file reports establishment counts by
county, NAICS, and employment size class, and publishes the three-way table
rather than only its margins.

CBP withholds cells for disclosure reasons. We measured 67.1\% of size-class
cells suppressed, and the rate climbs steadily with cell sparsity, from 9.9\% for
the smallest firm class to 98.5\% for the largest. A withheld cell cannot serve
as truth, so we carved out the largest dense, fully published sub-block:
\textbf{1,014 counties $\times$ 10 NAICS sectors $\times$ 4 size buckets}. That
is 40,560 atoms holding 6.1M establishments, 75.9\% of the national total, with
0.1\% zero cells and no imputation anywhere. The block is selected on
non-suppression, so it over-represents the well-populated part of the economy,
and every result below should be read as optimistic on that account.

We fit IPF under three margin regimes: one-way margins only, which is the
independence model; the two margins actually available in practice, county
$\times$ sector together with sector $\times$ size; and all three two-way
margins. Then we answer randomly generated queries and compare to truth. Errors
are symmetric relative error, $10^{|\log_{10}(\hat{y}/y)|} - 1$.

\section{Results}

\subsection{Individual cells are bad, aggregates are not}

Median per-atom error is 40.9\% under independence, 16.2\% under the realistic
two-margin regime, and 12.7\% with all three two-way margins. Individual cells
are unreliable no matter how the margins are chosen. Aggregation is what makes
the method work at all, and that shapes everything that follows.

Queries lying in $V$ came out exact on all 90 trials, at 0.00\%, which confirms
Section~\ref{sec:identity} numerically as well as algebraically. For box-shaped
queries, meaning the realistic market shape that constrains all three axes,
error falls steeply and monotonically with the share of total mass covered
(Table~\ref{tab:span}).

\begin{table}[h]
\centering
\begin{tabular}{lrrrr}
\toprule
Query mass share & $n$ & Median & p90 & p99 \\
\midrule
$<0.003\%$        & 262 & 11.39\% & 38.34\% & 95.49\% \\
$0.003$--$0.01\%$ & 344 &  6.77\% & 20.73\% & 41.34\% \\
$0.01$--$0.03\%$  & 412 &  4.39\% & 15.12\% & 25.88\% \\
$0.03$--$0.1\%$   & 631 &  2.99\% & 11.19\% & 29.88\% \\
$0.1$--$0.3\%$    & 651 &  2.34\% &  8.17\% & 16.15\% \\
$0.3$--$1\%$      & 729 &  1.52\% &  4.84\% & 11.65\% \\
$1$--$3\%$        & 564 &  0.73\% &  2.96\% &  7.23\% \\
$3$--$10\%$       & 353 &  0.39\% &  1.39\% &  3.51\% \\
$>10\%$           &  54 &  0.27\% &  0.67\% &  1.53\% \\
\bottomrule
\end{tabular}
\caption{Box-query error against query size, realistic margin regime. Median
error swings 42$\times$ across the range. Bin counts are given because the final
row rests on 54 queries, so its p99 is effectively a maximum and should not be
read as a quantile.}
\label{tab:span}
\end{table}

One result cuts against intuition. Uniformly random atom subsets were the
\emph{easiest} queries, at 0.27\% median error, because independent errors cancel
across many atoms. Box queries were consistently the hardest, since their
selection correlates with exactly the interaction structure IPF leaves out. Real
markets are boxes, so Table~\ref{tab:span} is the curve that matters, and the
scattered-query numbers are misleadingly reassuring.

\subsection{The statistic is calibrated, but does not out-rank query size}
\label{sec:calib}

We fit $\kappa$ in \eqref{eq:predictor} on 300 queries and test on 300 held out.
Our first comparison put it at Spearman $+0.663$ against $+0.522$ for a
query-size baseline, which looked like a decisive win. It was not. That baseline
had been handed an arbitrary constant instead of the same calibration
the new statistic received, and was scored on a slightly different subset. Run
fairly, with both predictors calibrated by the identical procedure on the
identical held-out set, the picture changes (Table~\ref{tab:predictor}).

\begin{table}[h]
\centering
\begin{tabular}{lrrrr}
\toprule
 & \multicolumn{2}{c}{All query families} & \multicolumn{2}{c}{Box queries only} \\
\cmidrule(lr){2-3}\cmidrule(lr){4-5}
Predictor & Spearman $\rho$ & $\mathrm{sd}(\log_{10})$ & Spearman $\rho$ & $\mathrm{sd}(\log_{10})$ \\
\midrule
$\lVert q_\perp\rVert / \langle q,H\rangle$ & $+0.721$ & 0.551 & $+0.602$ & 0.601 \\
query size alone                            & $+0.707$ & 0.915 & $+0.600$ & 0.810 \\
\bottomrule
\end{tabular}
\caption{Six seeds, 300 train and 300 test queries each, both predictors
calibrated identically and scored on the same subset. Standard deviations across
seeds are 0.045 and 0.042 for $\rho$.}
\label{tab:predictor}
\end{table}

On \emph{ranking}, the two are equivalent. The advantage of
$\lVert q_\perp \rVert$ averages $+0.013$ across all families and $+0.002$ on
box queries alone, both inside seed-to-seed noise. Anyone wanting only to sort
queries by riskiness can use query size and stop reading.

The advantage that survives is in \emph{magnitude}. Dispersion of prediction
error is consistently lower, $\mathrm{sd}(\log_{10})$ of 0.55 against 0.92 across
all families and 0.60 against 0.81 on boxes, holding on every seed tested. In
multiplicative terms, predictions land within roughly $3.5\times$ rather than
$6.5\times$. Calibration is also steadier: the identity-based statistic stayed in
$0.86$--$1.15\times$ across seeds while the query-size baseline wandered between
$0.34\times$ and $0.82\times$.

There is also one thing query size cannot do at any calibration. It cannot
return zero. Equation~\eqref{eq:predictor} is exactly zero for $q \in V$,
correctly identifying the queries that carry no error at all, and no function of
query size can express that.

For completeness: the normalised alignment ratio carries essentially no signal,
which is what sent us back to the unnormalised norm in the first place.

\subsection{Dollars cost a constant, not a regime change}

Establishment counts are the easiest thing to disaggregate, and a method that
only works on counts is not a market-sizing method. We tested two further
published sources at the same county $\times$ NAICS grain: Nonemployer Statistics
total receipts (2,863 counties $\times$ 12 sectors) and QCEW total annual wages
(668 counties $\times$ 12 sectors), holding lattice, protocol, and query
distribution fixed so only the measured quantity changes.

The penalty is a steady multiplicative factor near 1.5$\times$ in every
query-size bin: 1.34--1.63$\times$ for receipts and 1.23--2.22$\times$ for wages.
Nothing pathological happens. The mechanism shows up in dispersion, where the
within-county sector-share coefficient of variation rises from 0.39 for counts to
0.58 for receipts, and from 0.43 to 0.73 for wages. Since error falls sharply
with query size, a 1.5$\times$ penalty is bought back with roughly 3$\times$ more
span.

\subsection{How much modelling can add}

Is the residual modellable rather than merely bounded? Its spectrum says there is
something there: the leading component of the log-residual explains 24.3\% of
variance, against 3.6\% for a matched Gaussian benchmark. But it does not
generalise. Table~\ref{tab:ceiling} collects three attempts.

\begin{table}[h]
\centering
\begin{tabular}{llr}
\toprule
Method & Metric & Value \\
\midrule
Parametric interaction, transferred to held-out counties & median-error reduction & 9.4\% \\
Same form fitted in-sample (bound for that form)         & median-error reduction & 10.7\% \\
Low-rank completion (best at rank 1)                     & out-of-sample $R^2$    & 0.118 \\
County covariates (log size, sector diversity)           & out-of-sample $R^2$    & 0.041 \\
\bottomrule
\end{tabular}
\caption{Three approaches, two metrics. The transferred figure is a mean over five
train/test splits ($\mathrm{sd}\ 1.2$, range 8.2--11.5\%). The $R^2$ figures describe explained
variance of the log-residual and are not directly comparable to a median-error
reduction; converted to dispersion reduction they correspond to about 6.1\% and
2.1\%. All four therefore land in a single-digit band, 2\% to 11\%.}
\label{tab:ceiling}
\end{table}

We flag the metric distinction because it is easy to get wrong. Explained
variance and median-error reduction are different quantities, and reporting one
beside the other invites a comparison that is not valid. Once converted onto a
common footing, the four results agree closely, which is a good deal more
convincing than any single number would have been: roughly nine tenths of the
residual is cell-specific.

The transfer itself works well. Of whatever the parametric form can capture,
around 82\% survives being carried to held-out counties, so dimensionless structure does
move across geographies. There is simply not much of it to move. Low-rank
completion degrades to negative $R^2$ by rank 5, the usual signature of fitting
noise particular to the observed cells.

\subsection{Accumulated observations: local, and saturating}
\label{sec:storage}

If each past estimate is retained as a constraint, does the system improve? We
stored box observations and evaluated on held-out queries, first confirming that
stored constraints were satisfied to 0.000\% residual so a null result could not
be blamed on a bad fit.

Proximity decides everything. With 200 observations stored, error on unrelated
random markets moved from 2.57\% to 2.46\%; on markets neighbouring something
already stored it fell from 1.49\% to 0.16\%, with p90 dropping from 9.26\% to
3.69\%. Across five seeds the related gain is $87.2\% \pm 0.9$, about the most
stable figure in this paper, while the unrelated gain is $4.8\% \pm 3.2$ and went
negative on one seed. The second number is not distinguishable from zero.

Two further curves determine what a store is actually worth.
Table~\ref{tab:sat} varies how many observations sit near the query.

\begin{table}[h]
\centering
\begin{tabular}{lrrrrrrrr}
\toprule
Observations stored nearby & 0 & 1 & 2 & 3 & 5 & 8 & 13 & 21 \\
\midrule
Median error & 2.26\% & 0.24\% & 0.17\% & 0.13\% & 0.15\% & 0.30\% & 0.17\% & 0.16\% \\
Gain over none & n/a & 89\% & 93\% & 94\% & 93\% & 87\% & 93\% & 93\% \\
\bottomrule
\end{tabular}
\caption{The first nearby observation captures essentially all of the available
gain. The second and third add a sliver. Beyond that the series is flat and its
wobble is noise.}
\label{tab:sat}
\end{table}

Table~\ref{tab:decay} instead fixes eight stored observations and moves the
target away from them.

\begin{table}[h]
\centering
\begin{tabular}{lrrrrrrr}
\toprule
Atom overlap with stored region & 100\% & 95\% & 88\% & 75\% & 50\% & 25\% & 2\% \\
\midrule
Median error   & 0.09\% & 0.24\% & 0.62\% & 1.09\% & 1.06\% & 1.50\% & 1.54\% \\
Gain over none & 97\%   & 90\%   & 74\%   & 61\%   & 55\%   & 6\%    & none \\
\bottomrule
\end{tabular}
\caption{Benefit survives to roughly 50\% overlap and collapses between 50\% and
25\%.}
\label{tab:decay}
\end{table}

Together these say something sharper than locality alone. Because around 51\% of
a typical query already lies inside the margin span, a stored observation can
only speak to the remainder, and it does so within a neighbourhood of roughly
50\% overlap. Within that neighbourhood one observation is nearly sufficient.

\textbf{The value of such a store therefore scales with the number of distinct
neighbourhoods covered, not with the number of records held.} A thousand
observations of one market are worth about what three are worth. Three hundred
observations spread over three hundred neighbourhoods are worth roughly a
hundred times more than three hundred concentrated in one. This inverts the
intuition that a concentrated query stream is the favourable case: concentration
saturates, and diversity is the efficient direction per record. It also means
record count is the wrong health metric for such a system, and coverage is the
right one.

\section{Two failures that leave no trace}

\subsection{Margins that disagree}

IPF assumes its margins are true and mutually consistent. Margins pulled from
different sources, vintages, and definitions are neither. When they conflict, IPF
does not oscillate and does not diverge. It settles on a fixed point of the
composite map that satisfies whichever margin was applied \emph{last} and
approximates the other. Sweep-to-sweep drift reaches $10^{-16}$, which is
indistinguishable from clean convergence.

The practical effect is that loop order becomes an undocumented ranking of how
much you trust each source. With one margin at 2\% noise and another at 50\%, the
two orderings produced answers differing by 22.6\% of total mass, and both
reported convergence. Weighting each constraint by its precision removes the
order dependence and cuts median error from 41.7\% to 28.1\%, while costing
nothing when sources are equally good.

Margin noise on its own turns out to be second-order: 10\% multiplicative noise
moved median cell error only from 16.2\% to 21.3\%. The damage comes from
treating unequal sources as equal, not from noise as such. Where noisy estimates
exist at several aggregation levels at once, IPF is the wrong tool entirely and
optimal reconciliation is the right one, since it projects incoherent forecasts
onto the coherent subspace using the full error covariance and handles non-nested
groupings~\cite{mint}.

\subsection{Suppression is adversarial by design}

The 67\% suppression we ran into is not incidental thinness in the data.
Agencies apply \emph{primary} suppression to sensitive cells, then
\emph{complementary} suppression to further cells chosen specifically so the
primary values cannot be reconstructed from the published
margins~\cite{cellsuppression}. Inference from margins is the stated threat
model. This is an information barrier rather than a modelling problem, and it
falls hardest on small, specific queries, which Table~\ref{tab:span} already
identifies as the least reliable regime. The two effects compound.

\section{Discussion}

\paragraph{A rule you can operate.} Section~\ref{sec:calib} means the operating
rule can be stated on query size alone, which is the simpler instrument; the
identity buys a tighter magnitude and the ability to certify exact queries, not
a better ordering. Above roughly 1\% of the anchor's mass,
median error runs below 1\% and p90 below 3\%. Below roughly 0.01\%, p90 passes
15\% and p99 approaches 100\%. Since \eqref{eq:predictor} can be evaluated before
answering, the honest response in the second regime is to decline rather than to
print a number with decimals on it. For dollar-denominated queries the thresholds
move by about 3$\times$ and $\kappa$ needs its own calibration.

\paragraph{Where the leverage actually is.} Ranking every intervention we
measured: observing an \emph{uncovered} cell cuts error by 87\%; combining
sources by precision, through a Fay--Herriot shrinkage estimator, by 33.5\%;
better structural modelling by 2\% to 11\%. The qualifier matters, because
Table~\ref{tab:sat} shows a second observation of an already-covered cell buys
almost nothing. Effort belongs in acquisition and in fusing
evidence, not in a smarter interpolator. One implementation note, learned by
getting it wrong: because market quantities span orders of magnitude, shrinkage
has to run on a scale where variances are roughly homoskedastic. Fitted on the
raw scale it was \emph{worse} than either of its inputs at high noise. In log
space it beat both at every noise level we tested.

\paragraph{Coherence is not accuracy.} Reconciliation forces internal consistency
whether or not the result is true. Any system that catches defects by spotting
contradictions loses that signal the moment a reconciliation layer is added, and
has to replace it with accuracy checks. The held-out marginal protocol used here
is one such check, and it can be run continuously rather than once.

\paragraph{What this does not cover.} Our ground truth is a dense sub-block
selected on non-suppression, which biases everything optimistic. The lattice is
coarse at 10 sectors and 4 size bands, far coarser than a real market definition
needs. Receipts and wages are proxies for spend, and neither separates price from
frequency. Establishment populations are not the buyer population for most
products, and the lattice indexes \emph{sellers}. For a local business the
seller's geography and the buyer's coincide, which conceals the distinction; for
a business selling online they are unrelated, and a pipeline that silently uses
the seller's county as the market's geography will produce a well-formed answer
to the wrong question. Our queries are boxes with binary membership, whereas real market
definitions carry fractional weights that are themselves estimated and probably
contribute more error than anything we measured. And none of this touches market
definition: whether two products compete is not a statistical question, and no
amount of lattice resolution will answer it.

\section{Conclusion}

Disaggregating a published total is a sound procedure whose reliability varies by
more than two orders of magnitude across the queries people actually ask, with
nothing in the output to distinguish the good cases from the bad. The residual of
a margin-constrained fit is orthogonal to what was measured, which makes the
reliability of any particular query computable in advance. That statistic will
not rank queries better than their size does, but it is better calibrated, it
identifies exactly answerable queries, and it can be evaluated before committing
to an answer.

The negative results are the more useful half. Modelling the residual is capped
near a tenth of it. Accumulating observations helps only nearby and saturates
after one. Suppression is engineered against precisely this kind of inference.
Taken together they point away from better estimators and toward acquiring
observations in uncovered neighbourhoods and fusing sources by measured
precision, which is roughly what the federal statistical system already does.

\section{Related work}

IPF is due to Deming and Stephan~\cite{deming1940}, with convergence to the
$I$-projection established by Csisz\'ar~\cite{csiszar1975} and the definitive
treatment in R\"uschendorf~\cite{ruschendorf1995}. The algorithm coincides with
Sinkhorn--Knopp, and the optimal-transport literature supplies sharp geometric
convergence rates~\cite{sinkhornrate} along with the log-domain stabilisation
needed when margins approach zero. Decomposable constraint sets admit closed-form
solutions~\cite{decomposable}.

Applied work on spatial microsimulation reports the same sensitivity to sparsity
we observe, and notes that a quality measure for the resulting estimates is
generally unavailable~\cite{jasss2015}. That literature validates against
aggregate constraint totals rather than a held-out joint, which is the one place
our protocol improves on it. Estimation for domains too small to measure directly
is the province of small-area estimation, from Fay and
Herriot~\cite{fayherriot1979} through the survey of Rao and
Molina~\cite{rao2015}; it supplies principled mean-squared-error estimates and
spatial borrowing of strength, and subsumes much of what we build here.
Reconciling incoherent hierarchical forecasts is treated by Wickramasuriya et
al.~\cite{mint}. When constraints are both noisy and insufficient, generalised
maximum entropy is the natural setting~\cite{golan1996}.

\begin{thebibliography}{99}
\bibitem{deming1940} W.~E. Deming and F.~F. Stephan. On a least squares
adjustment of a sampled frequency table when the expected marginal totals are
known. \emph{Ann. Math. Statist.}, 11(4):427--444, 1940.
\bibitem{csiszar1975} I.~Csisz\'ar. $I$-divergence geometry of probability
distributions and minimization problems. \emph{Ann. Probab.}, 3(1):146--158, 1975.
\bibitem{ruschendorf1995} L.~R\"uschendorf. Convergence of the iterative
proportional fitting procedure. \emph{Ann. Statist.}, 23(4):1160--1174, 1995.
\bibitem{sinkhornrate} On the convergence rate of Sinkhorn's algorithm.
arXiv:2212.06000.
\bibitem{decomposable} Iterative proportional scaling via decomposable submodels
for contingency tables. arXiv:math/0603495.
\bibitem{jasss2015} Evaluating the performance of iterative proportional fitting
for spatial microsimulation. \emph{J. Artificial Societies and Social
Simulation}, 18(2):21, 2015.
\bibitem{fayherriot1979} R.~E. Fay and R.~A. Herriot. Estimates of income for
small places: an application of James--Stein procedures to census data.
\emph{JASA}, 74(366):269--277, 1979.
\bibitem{rao2015} J.~N.~K. Rao and I.~Molina. \emph{Small Area Estimation}.
Wiley, 2nd edition, 2015.
\bibitem{mint} S.~L. Wickramasuriya, G.~Athanasopoulos, and R.~J. Hyndman.
Optimal forecast reconciliation for hierarchical and grouped time series through
trace minimization. \emph{JASA}, 114(526):804--819, 2019.
\bibitem{cellsuppression} R.~Zhang, L.~Chen, and Y.~Cheng. Overview of cell
suppression methods. \emph{Proc. ASA Survey Research Methods Section}, 2023.
\bibitem{golan1996} A.~Golan, G.~G. Judge, and D.~Miller. \emph{Maximum Entropy
Econometrics: Robust Estimation with Limited Data}. Wiley, 1996.
\end{thebibliography}

\end{document}
